\chapter{Wprowadzenie}\label{ch:wprowadzenie}

\section{Tło techniczne}\label{sec:tlo_techniczne}

Współczesne złośliwe oprogramowanie, takie jak trojany zdalnego dostępu (RAT) oraz programy wykradające dane (Infostealer), powszechnie wykorzystuje techniki omijania klasycznych mechanizmów obronnych. Zamiast zapisywać pliki wykonywalne bezpośrednio na dysku, atakujący opierają wykonanie na wbudowanych narzędziach systemowych (\textit{Living off the Land Binaries}, LOLBins) oraz bezpośredniej manipulacji pamięcią operacyjną procesów. Do najczęściej spotykanych technik należą:
\begin{itemize}
    \item \textbf{Wstrzykiwanie kodu (\textit{Process Injection}):} alokacja pamięci wirtualnej w przestrzeni adresowej innego, działającego procesu i zapisanie w niej złośliwego ładunku wraz z uruchomieniem zdalnego wątku.
    \item \textbf{Drążenie procesów (\textit{Process Hollowing}):} utworzenie legalnego procesu systemowego (np. \code{svchost.exe}, \code{explorer.exe}) w stanie zawieszonym, usunięcie jego oryginalnej sekcji kodu z pamięci i zastąpienie jej obrazem złośliwego programu przed wznowieniem wykonania.
    \item \textbf{Refleksyjne ładowanie bibliotek (\textit{Reflective DLL Injection}):} ładowanie biblioteki DLL bezpośrednio z bufora pamięci RAM przez własny kod inicjalizujący, z pominięciem standardowego menedżera ładowania systemu Windows i bez rejestracji modułu na liście załadowanych bibliotek procesu.
\end{itemize}

W rezultacie programy antywirusowe oparte na statycznych sygnaturach plikowych często nie wykrywają zagrożenia. Głównym źródłem danych o infekcji staje się pamięć RAM, w której znajdują się rozpakowane ładunki binarne i konfiguracje, oraz szczegółowe dzienniki zdarzeń systemowych rejestrujące zachowanie procesów w czasie rzeczywistym.

\section{Problem badawczy}\label{sec:problem_badawczy}

Zespoły zajmujące się monitorowaniem bezpieczeństwa i reagowaniem na incydenty (SOC/DFIR) mierzą się z problemem niskiej trwałości prostych wskaźników kompromitacji (IoC). Zgodnie z modelem Piramidy Bólu Davida Bianco \cite{bianco2013pyramid}, sumy kontrolne plików czy adresy IP i domeny serwerów C2 są bardzo łatwe do zmiany przez atakującego poprzez ponowne spakowanie pliku lub rotację infrastruktury sieciowej.

Trwałą skuteczność detekcji daje dopiero identyfikacja wskaźników aktywności (IoA) oraz wzorców zachowań (TTP) sklasyfikowanych w bazie MITRE ATT\&CK \cite{mitre_attack}. Opracowanie takich reguł wymaga jednak połączenia dwóch uzupełniających się perspektyw:
\begin{enumerate}
    \item analizy zrzutu pamięci RAM, która pokazuje stan procesów, wstrzyknięte sekcje kodu oraz niezaszyfrowane ciągi znaków w momencie infekcji,
    \item analizy dzienników zdarzeń systemowych i ruchu sieciowego, które pozwalają odtworzyć chronologiczny przebieg ataku od uruchomienia próbki po utrwalenie obecności w systemie.
\end{enumerate}

Zestawienie obu tych źródeł pozwala tworzyć zarówno reguły wykrywające ładunki bezpośrednio w pamięci operacyjnej (YARA \cite{yara_docs}), jak i reguły wykrywające podejrzane sekwencje zdarzeń w logach systemowych (Sigma \cite{sigma_hq}).

\section{Cel i zakres projektu}\label{sec:cele_projektu}

Głównym celem projektu jest zbudowanie odizolowanego środowiska laboratoryjnego oraz przeprowadzenie poinfekcyjnej analizy pamięci RAM i dzienników zdarzeń systemu Windows dla wybranych próbek złośliwego oprogramowania, zakończonej opracowaniem i weryfikacją reguł detekcyjnych YARA i Sigma oraz zestawieniem wskaźników IoC/IoA.

Cele szczegółowe obejmują:
\begin{enumerate}
    \item konfigurację bezpiecznego środowiska wirtualnego z symulacją usług sieciowych i rozszerzonym logowaniem zdarzeń systemowych,
    \item kontrolowane uruchomienie badanych próbek złośliwego oprogramowania oraz pobranie materiału dowodowego (zrzutu pamięci RAM, dzienników zdarzeń i ruchu sieciowego),
    \item przeprowadzenie analizy pamięci operacyjnej oraz korelacji logów w celu rekonstrukcji przebiegu infekcji i wyodrębnienia wstrzykniętych ładunków,
    \item przygotowanie i przetestowanie autorskich reguł detekcyjnych w formatach YARA oraz Sigma wraz ze zbiorczym zestawieniem wykrytych wskaźników kompromitacji i aktywności.
\end{enumerate}

\subsection*{Zakres i założenia projektowe}
\begin{itemize}
    \item \textbf{System operacyjny:} Microsoft Windows 10 x64.
    \item \textbf{Izolacja sieciowa:} środowisko jest całkowicie odcięte od sieci zewnętrznej (WAN), a komunikacja sieciowa złośliwego oprogramowania odbywa się wyłącznie z lokalnym serwerem symulującym usługi internetowe.
    \item \textbf{Faza analizy:} projekt koncentruje się na analizie poinfekcyjnej - badane jest działanie próbki po jej uruchomieniu w systemie, z pominięciem sposobu dostarczenia do ofiary.
\end{itemize}
